\section{Greedy与ILS}
\label{sec:03_03}

贪心构造首先解决如何从空选择得到一个可行方案；迭代局部搜索则允许撤销部分选择，在新的组合中继续寻找更高收益。两者共同作用于同一个带容量的超图边选择问题，不改变任务收益或资源约束。本节依据Großmann等的贪心与局部搜索框架\cite{GrossmannJRS2026}，说明已有诊断和压力实验实际调用的构造规则、交换邻域、扰动、接受条件及停止方式。这里使用的贪心优先级在接口中名为\texttt{bmindegree\_dynamic}，与原作者完整流程采用的$S,\mathrm{pin}$配置不同；本节调用也不包含完整的归约与解恢复过程。

\subsection{可行选择与贪心优先级}

设$H=(V,E)$为输入超图，$V$是资源集合，$E$是候选任务索引集合，$m=|V|$与$n=|E|$分别表示资源数和任务数。每个任务$e\in E$需要非空资源集合$R_e\subseteq V$，具有无量纲正收益$w_e$。一般数学模型允许$w_e\in\mathbb{R}_{>0}$；本节实际适配器使用整数收益，因此实现层面的输入满足$w_e\in\mathbb{Z}_{>0}$。资源$v\in V$的容量$b_v\in\mathbb{Z}_{>0}$表示同时允许选择的关联任务数，每个任务在其关联资源上占用一个容量单位。容量、占用次数和度均为无量纲计数。

对任务集合$M\subseteq E$，定义总收益和资源负载为
\begin{align}
 W(M)&=\sum_{e\in M}w_e,
 \label{s0303:weight}\\
 \ell_M(v)&=\sum_{e\in M}\mathbb{I}_{\{v\in R_e\}},
 &&v\in V.
 \label{s0303:load}
\end{align}
其中$\mathbb{I}_{\{P\}}$为命题$P$的指标函数：当$P$成立时取$1$，否则取$0$。集合$M$可行当且仅当$\ell_M(v)\leq b_v$对全部$v\in V$成立。对可行$M$，剩余容量$c_M(v)=b_v-\ell_M(v)$属于$\{0,\ldots,b_v\}$。当前可以直接加入的任务构成集合
\begin{equation}
 \mathcal{F}(M)=\left\{e\in E\setminus M:
 c_M(v)\geq1,\ \forall v\in R_e\right\}.
 \label{s0303:free_edges}
\end{equation}
该条件同时检查一个任务需要的全部资源，不能以平均剩余容量为正替代。

记输入图上资源$v$的度为
\begin{equation}
 d_H(v)=\left|\{e\in E:v\in R_e\}\right|,
 \qquad v\in V.
 \label{s0303:degree}
\end{equation}
它统计传入当前构造过程的图中使用该资源的候选任务数，不随当前选择$M$改变。孤立资源的度可以为零，但只要$v\in R_e$，便有$d_H(v)\geq1$，因此下式的分母非零。实际贪心规则对应的理想实数优先级为
\begin{equation}
 h_M(e)=w_e\prod_{v\in R_e}
 \frac{b_v-\ell_M(v)}{d_H(v)},
 \qquad e\in\mathcal{F}(M).
 \label{s0303:priority}
\end{equation}
$h_M(e)>0$为无量纲优先级。较高收益提高优先级；较多剩余容量意味着选择该任务后资源仍可能服务其他任务；较大的资源度则说明更多候选任务竞争该资源。乘积将一个任务涉及的各资源共同纳入排序，但它不是选择该任务所造成的真实收益损失，也不构成目标函数的上界。

式\eqref{s0303:priority}中的$d_H(v)$不是剩余可行任务中使用$v$的任务数。若将其替换为随$M$变化的可行邻居数，便改变了本批实验的算法。对于全部已有组合实验采用的单位容量$b_v=1$，任一$e\in\mathcal{F}(M)$都满足$\ell_M(v)=0$，$v\in R_e$，于是
\begin{equation}
 h_M(e)=\frac{w_e}{\displaystyle\prod_{v\in R_e}d_H(v)},
 \qquad e\in\mathcal{F}(M),\quad b_v=1\ (v\in V).
 \label{s0303:unit_priority}
\end{equation}
因此，只要一个候选仍然可行，其理想分数在这一单位容量构造过程中保持不变。每次加入任务后发生变化的是候选的可行状态，不能从接口名称中的动态二字推断每一步都会改变所有候选分数。

上述公式描述理想实数运算。实际适配器使用\texttt{double}，按资源关联项逐项累乘分子和度乘积，最后相除，并未采用对数求积。有限精度下，乘积溢出或商下溢可能使数学上为正的优先级成为机器零值，从而增加排序中的并列。对保留的\texttt{us04}输入进行独立表达式核验时，在输入哈希核对一致的条件下，按源码相同的双精度乘积重算其$163$个任务的初始优先级，发现其中$118$个任务的度乘积溢出为正无穷，所得机器优先级为$0$，而其理想实数优先级仍为正。这一核验只检查初始优先级表达式，没有重跑贪心或ILS，也不能据此认定某项历史收益或耗时由该数值现象造成。历史实现及正式成绩均原样保留；正无穷乘积不是新增的数学排序规则，本书也不据此修改公开算法。第3.8节进一步讨论这一实现边界的适用范围。

贪心从$M=\varnothing$开始，按实际计算的优先级由高到低检查候选，通过堆组织排序，并根据当前资源状态过滤不可行任务。在理想运算下，单位容量构造可理解为按式\eqref{s0303:unit_priority}组织候选顺序，再逐项执行式\eqref{s0303:free_edges}的检查；实际顺序还受上述数值并列及其处理方式影响。接收任务$e$时更新
\begin{equation}
 M\leftarrow M\cup\{e\},\qquad
 \ell_M(v)\leftarrow\ell_M(v)+1,
 \quad v\in R_e,
 \label{s0303:greedy_add}
\end{equation}
其余资源负载不变。这里负载赋值的右端表示接收前的负载。由于构造阶段只增加任务，已因容量不足被过滤的候选不会在后续重新变得可行。

构造结束后，$\mathcal{F}(M)=\varnothing$，即不存在能够直接追加的任务。这样的集合称为极大可行选择。极大性仅排除了追加一个任务的改进，并不排除撤销已有任务后获得更高收益；最大收益解则必须在全部可行集合中达到最高收益。这一区别正是引入交换和扰动的原因。

\subsection{删一添二交换与极大化补全}

局部搜索检查能否删除一个已选任务，换入两个未选任务。设$e\in M$为拟删除的任务，$f,g\in E\setminus M$为两个不同的拟加入任务，$f\ne g$。交换后的集合为
\begin{equation}
 M'=(M\setminus\{e\})\cup\{f,g\}.
 \label{s0303:swap_set}
\end{equation}
这一操作通常称为$(1,2)$交换，其中第一个数字表示删除数量，第二个数字表示加入数量。交换的精确容量条件是
\begin{equation}
 \ell_M(v)-\mathbb{I}_{\{v\in R_e\}}
 +\mathbb{I}_{\{v\in R_f\}}
 +\mathbb{I}_{\{v\in R_g\}}
 \leq b_v,
 \qquad\forall v\in V.
 \label{s0303:swap_capacity}
\end{equation}
左端逐项扣除被删除任务的占用，再加入两个新任务的占用。分别验证$f$和$g$能否单独加入还不够，因为二者可能争用同一份剩余容量。收益严格改进还要求
\begin{equation}
 \Delta W=w_f+w_g-w_e>0.
 \label{s0303:swap_gain}
\end{equation}
$\Delta W$与$W(M)$均为无量纲收益。等收益交换不属于这里的改进交换。

在本批实验的单位容量条件下，可以用阻挡集合更直接地描述候选筛选。对未选任务$f$，定义
\begin{equation}
 B_M(f)=\{a\in M:R_a\cap R_f\ne\varnothing\},
 \qquad f\in E\setminus M.
 \label{s0303:blockers}
\end{equation}
$B_M(f)$列出所有与$f$争用资源的已选任务。若$M$已经极大，则每个未选任务都至少有一个阻挡者。删除$e$后刚好能够加入的邻接候选属于
\begin{equation}
 \mathcal{N}_1(e;M)=
 \{f\in E\setminus M:B_M(f)=\{e\}\},
 \qquad e\in M.
 \label{s0303:one_blocker}
\end{equation}
下标$1$表示只有一个阻挡任务，不表示只共享一个资源；$e$和$f$可以共享多个资源，但删除同一个$e$便可同时释放这些资源。单位容量下，从$\mathcal{N}_1(e;M)$中选出不同的$f,g$后，交换可行性与收益条件化为
\begin{equation}
 R_f\cap R_g=\varnothing,
 \qquad w_f+w_g>w_e.
 \label{s0303:unit_swap}
\end{equation}
前一条件保证新任务之间也不冲突，后一条件保证替换后收益严格上升。一般容量下，两个任务共享资源仍可能可行，因而必须回到式\eqref{s0303:swap_capacity}，不能直接沿用无交集条件。

公开内核围绕已选任务收集受阻邻接候选，并利用阻挡计数和扫描标记减少重复检查；候选列表按收益降序排列，再查找满足条件的任务对。这里对实际交换行为的说明以本批单位容量输入为范围，式\eqref{s0303:swap_capacity}则给出一般容量下应满足的数学可行性条件。每次成功交换后，还执行极大化补全：只要$\mathcal{F}(M')$非空，就从其中取出一个任务加入，直至没有可直接加入的任务。记这一过程的输出为$\mathcal{A}(M')$，则
\begin{equation}
 M'\subseteq\mathcal{A}(M'),\qquad
 \mathcal{F}(\mathcal{A}(M'))=\varnothing,\qquad
 W(\mathcal{A}(M'))\geq W(M').
 \label{s0303:maximize}
\end{equation}
$\mathcal{A}$表示按实现中的遍历顺序补全，不表示求解一个新的全局最大化问题；不同补全顺序可能得到不同的极大集合。

重复交换与补全，直到一轮检查不再找到严格改进。以$\mathcal{L}(M)$表示这一穷尽改进过程的输出。每次成功交换都使收益严格增加，补全加入的任务也具有正收益，而有限输入只有有限个任务子集，因此严格改进序列不会无限延长。其停止含义限于所使用的交换搜索：它不提供全局最优证书，也不能排除需要同时删除更多任务才能实现的改进。

\subsection{扰动如何改变局部搜索的起点}

单纯重复严格改进交换会停留在当前交换规则无法继续改善的状态。扰动通过尝试强制加入未选任务，并移除阻止其加入的已选任务，构造另一个可行起点。设单位容量下拟强制加入的任务为$q\in E\setminus M$，则一次成功强制插入在集合层面对应
\begin{equation}
 \widetilde M=(M\setminus B_M(q))\cup\{q\}.
 \label{s0303:forced_insert}
\end{equation}
删除$B_M(q)$释放$q$需要的全部资源，而其他已选任务之间的可行性不变。因此$\widetilde M$可行，但其收益变化为
\begin{equation}
 W(\widetilde M)-W(M)
 =w_q-\sum_{a\in B_M(q)}w_a,
 \label{s0303:perturb_gain}
\end{equation}
该差值可以为负。扰动的作用是改变后续局部搜索能够到达的组合，而不是保证每次插入立即获益。

实际内核首先根据当前匹配的任务数量$|M|$生成整数$\alpha\in\{1,\ldots,2|M|\}$，初始化扰动长度$k=1$。当$\alpha=1$时，另设$i=1$，每当硬币调用返回$1$就令$i\leftarrow i+1$，在首次返回其他结果后令$k=i+1$。因此，普通分支尝试一个候选，另一分支产生至少为$2$的扰动长度。这里$|M|$、$\alpha$、$i$和$k$均为无量纲整数；它们描述当前方案大小和扰动操作数量，与目标收益$W(M)$不同。该调用要求当前匹配非空，适配器通过贪心收益为正的进入条件排除空起点。

当$k=1$时，内核调用其未选受阻候选选择函数获取任务$q$，逐一检查$q$的关联资源。若某资源已经满载，则按关联任务的遍历次序寻找一个允许移除的已选任务，释放一个容量单位。该分支还检查源码中的精确标记，以避免删除受保护任务。最终仅在$q$确实可加入时执行插入，因而一次尝试不能无条件等同于一次成功替换。

当$k>1$时，内核先调用候选选择函数四次，再从四个返回候选中选取内部年龄字段数值最小的任务$x$作为起点。这四次调用的返回值不必互异，年龄字段也不是以秒计量的物理时间。随后，代码沿资源关联建立一个候选列表：从$p_1\in R_x$出发，经使用$p_1$的任务$e_1$，再经$p_2\in R_{e_1}$且$p_2\ne p_1$，收集使用$p_2$且不同于$e_1$的任务$e_2$。最后把$x$也加入列表。其中$p_1,p_2\in V$为资源索引，$x,e_1,e_2\in E$为任务索引；这条关联路径用于生成扰动候选，不表示任务执行的先后顺序。

该结构是列表而非去重后的集合，同一任务可由不同关联路径多次进入列表。列表随机打乱后，实际处理的条目数为
\begin{equation}
 k'=\min\{k,L\},
 \qquad k'\in\mathbb{Z}_{\geq0},
 \label{s0303:perturb_length}
\end{equation}
其中$L$为列表长度，是无量纲计数。代码依次处理前$k'$个条目：对仍未选中的任务，释放其满载资源上的占用，再将其加入；对当时已经选中的任务则跳过。先前插入会改变后续条目的阻挡关系，所以该过程不能解释为同时加入$k'$个互不影响的任务，也不能把$k'$视为最终新增任务数。

两个分支结束后均执行极大化补全。用$\mathcal{P}(M)$表示扰动并补全后的可行集合，则一轮迭代局部搜索产生的候选解为
\begin{equation}
 M^{\mathrm{new}}=\mathcal{L}(\mathcal{P}(M)).
 \label{s0303:ils_candidate}
\end{equation}
这一表达式固定了操作次序：先通过扰动改变组合，再补全并寻找严格改进交换。候选选择函数、容器遍历和随机生成器的具体行为属于源码复现的一部分；仅凭这里的关联关系，不能将扰动替换为从全部未选任务中独立均匀抽样。

\subsection{当前解、历史最好解与离散接受规则}

迭代局部搜索同时维护当前解$M$与历史最好解$M^{\mathrm{best}}$。当前解决定下一次扰动的起点，历史最好解用于保存搜索过程中已经取得的最高收益。二者初始均等于贪心输出$M^{\mathrm{G}}$：
\begin{equation}
 M=M^{\mathrm{best}}=M^{\mathrm{G}}.
 \label{s0303:initial_states}
\end{equation}
实际内核从这一状态直接进入扰动，再进行交换改进；不能在复现时无说明地增加一次扰动前的独立局部搜索。

若$W(M^{\mathrm{new}})>W(M)$，候选自动成为新的当前解。否则，在中间整数计算无溢出时，源码的离散接受规则可以写为
\begin{align}
 D={}&1+\bigl(W(M)-W(M^{\mathrm{new}})\bigr)
          \bigl(W(M^{\mathrm{best}})-W(M^{\mathrm{new}})\bigr),
 \label{s0303:accept_denominator}\\
 U\sim{}&\operatorname{Unif}\{1,\ldots,D\}.
 \label{s0303:accept_draw}
\end{align}
$D\in\mathbb{Z}_{\geq1}$为离散抽样区间的上端点，$U$为在该整数区间均匀抽取的随机数。这里使用了实际输入的整数收益，以及历史最好收益不低于当前收益这一不变量。仅当$U=1$时接受候选，因此非改善候选的接受概率为
\begin{equation}
 p_{\mathrm{acc}}=
 \frac{1}{1+\bigl(W(M)-W(M^{\mathrm{new}})\bigr)
                \bigl(W(M^{\mathrm{best}})-W(M^{\mathrm{new}})\bigr)}.
 \label{s0303:accept_probability}
\end{equation}
全部收益均为无量纲数，故分母与概率也无量纲。若候选与当前解等收益，第一个差值为零，$D=1$，候选必然被接受；若候选更差，则两个差值均为正，接受概率小于$1$，并随收益差的增大而降低。这是允许较差当前解的明确规则，并不存在一个单独给定的稍劣幅度阈值。上述概率解释以中间整数计算无溢出为条件，并不是对任意输入规模下机器整数运算安全性的证明；现有材料也不据此断言已发生整数溢出。

论文给出的概率表达式没有式\eqref{s0303:accept_probability}分母中的$1$，而本批调用的公开内核采用上述整数抽样。该差别尤其影响等收益边界：实际代码在此时有确定的接受行为。复现实验应以代码的离散规则为准，不能直接用论文中省略$1$的表达式替代。

接受或拒绝候选后，只有当当前解严格超过历史最好收益时，才更新$M^{\mathrm{best}}$。因此，允许当前解退步不会覆盖已经保存的最好解。若以$t\in\mathbb{Z}_{\geq0}$表示完成的外层迭代数，则在各轮正常执行、收益计算与比较正确的条件下，有
\begin{equation}
 W(M^{\mathrm{best}}_{t+1})\geq W(M^{\mathrm{best}}_t),
 \qquad
 W(M^{\mathrm{best}}_t)\geq W(M_t).
 \label{s0303:best_invariants}
\end{equation}
前一个关系表示最好收益单调不降，后一个关系区分保存结果与当前搜索位置。最终返回的是$M^{\mathrm{best}}$，而不一定是最后一轮的当前解。

\subsection{停止计数与完整数学流程}

设$K\in\mathbb{Z}_{\geq0}$为连续未刷新历史最好收益的停止阈值，$q\in\mathbb{Z}_{\geq0}$为相应计数器；这里$q$表示计数，不再表示扰动中的候选任务。源码中的计数顺序是先令$q\leftarrow q+1$，再检查是否创造新的历史最好收益；只有严格刷新最好收益时，才令$q\leftarrow0$。随后检查
\begin{equation}
 q>K.
 \label{s0303:stop}
\end{equation}
所以，连续$K$轮未刷新纪录并不会立即停止，在没有其他停止条件时，要到第$K+1$轮连续未刷新纪录后才停止。若当前解先因接受较差候选而下降，之后有所回升，但仍未超过历史最好收益，计数也不会重置。

在不启用时间限制的情况下，本节适配器对应的过程如下。
\begin{enumerate}
 \item 根据资源关联和容量构造输入图，从空集执行贪心，得到可行极大集合$M^{\mathrm{G}}$。
 \item 若接口参数为$-1$，直接返回$M^{\mathrm{G}}$，不调用ILS。若贪心收益为零，适配器也不进入ILS。
 \item 对$K\geq0$且$W(M^{\mathrm{G}})>0$的情况，令$M=M^{\mathrm{best}}=M^{\mathrm{G}}$，$q=0$。
 \item 从当前解产生$M^{\mathrm{new}}=\mathcal{L}(\mathcal{P}(M))$。若其收益严格高于当前收益，令$M\leftarrow M^{\mathrm{new}}$；否则按式\eqref{s0303:accept_denominator}和式\eqref{s0303:accept_draw}抽取$U$，仅当$U=1$时执行同一赋值。
 \item 令$q\leftarrow q+1$。若$W(M)>W(M^{\mathrm{best}})$，令$M^{\mathrm{best}}\leftarrow M$并令$q\leftarrow0$。
 \item 若$q>K$，返回$M^{\mathrm{best}}$；否则回到第4步。
\end{enumerate}

其中$-1$是适配器用于跳过ILS的控制值，不是连续未改进次数的数学取值。原公开函数还提供默认值为$1800\,\mathrm{s}$的可选时间上限，但本节适配器显式传入无穷大，将这一时间停止条件关闭。因此不能把这些运行解释为统一的$1800\,\mathrm{s}$实验。

$K=15$、$1000$和$5000$均表示连续未改进计数阈值，不是固定的外层迭代总数。每次刷新历史最好收益都会重新计数，一次外层迭代内部还可能执行多次交换与补全。由此，阈值相同并不意味着不同实例执行相同数量的交换，也不意味着消耗相同的运行时间。对有限输入，严格刷新历史最好收益的次数有限；在每次扰动和局部搜索调用均正常结束、计数及收益运算正确的前提下，有限$K$最终会触发式\eqref{s0303:stop}，但这不给出与$K$相等的总迭代预算。

由于最好解独立保存，在正常完成、保持可行性且收益计算正确的调用中，返回值满足
\begin{equation}
 W(M^{\mathrm{best}})\geq W(M^{\mathrm{G}}).
 \label{s0303:not_worse_than_greedy}
\end{equation}
该关系只保证同一次调用不丢失其贪心起点，既不保证严格改善，也不保证达到全局最优。它还不能推出两个独立运行中较大的$K$必然产生更高收益，因为二者未必经过相同的搜索轨迹。

\subsection{已有调用范围与比较口径}

已有诊断实验在三个条件上分别使用接口参数$-1$、$1000$和$5000$，共形成九次运行。其中$-1$对应纯贪心，另外两种配置分别启用上述连续未改进阈值。诊断调用使用随机种子$0$，但固定这一种子只控制相应随机生成器，不能保证散列表遍历、同分候选次序及全部执行路径完全相同。因此，两个阈值的独立运行不能直接解释为同一条随机搜索轨迹的短前缀和长前缀；其质量差异须根据实际记录判断。

压力实验在$59$个条件上采用Greedy加ILS，阈值为$1000$。这些调用沿用本节的贪心配置，不包含原作者完整的归约、$S,\mathrm{pin}$构造和解恢复流程。完整流程中的ILS阈值$15$与这里的$1000$、$5000$具有相同的计数含义，但整体算法组成和运行条件不同，不能仅凭阈值大小推断总预算或性能高低。

适配器将内部时间分为建图、贪心和ILS三个阶段。设四个单调时钟读数$t_0,t_1,t_2,t_3$依次位于建图准备开始、匹配结构建立完成、贪心结束和可选ILS调用结束，则
\begin{align}
 T_{\mathrm{build}}&=t_1-t_0,\label{s0303:build_timing}\\
 T_{\mathrm{greedy}}&=t_2-t_1,\label{s0303:greedy_timing}\\
 T_{\mathrm{ILS}}&=t_3-t_2.
 \label{s0303:timing}
\end{align}
三个时间均非负，单位为$\mathrm{s}$。建图阶段包含随机种子设置、任务资源列表建立、必要的列表排序、容量赋值及匹配结构初始化；贪心阶段记录构造调用，ILS阶段记录可选局部搜索调用所在的时间区间。参数为$-1$时，$T_{\mathrm{ILS}}$对应的区间仍会被计时，但没有执行ILS，不能将其中极小的程序开销解释为搜索耗时。已选任务标识在$t_3$之后复制到输出数组，这部分不计入上述三个内部阶段。

这些定义允许分别观察构造与改进的开销，却不能消除不同实验批次在停止条件、预处理、热启动和CPU并发上的差异。第3.7节据实际逐次记录比较收益与耗时；本节的算法性质不预先排除未改善、未达最优或更大阈值未取得更高收益的结果。
